%%
%% This is file sample-acmcp.tex',
%% generated with the docstrip utility.
%%
%% The original source files were:
%%
%% samples.dtx  (with options: all,journal,acmcp')
%% 
%% IMPORTANT NOTICE:
%% 
%% For the copyright see the source file.
%% 
%% Any modified versions of this file must be renamed
%% with new filenames distinct from sample-acmcp.tex.
%% 
%% For distribution of the original source see the terms
%% for copying and modification in the file samples.dtx.
%% 
%% This generated file may be distributed as long as the
%% original source files, as listed above, are part of the
%% same distribution. (The sources need not necessarily be
%% in the same archive or directory.)
%%
%%
%% Commands for TeXCount
%TC:macro~\cite [option:text,text]
%TC:macro~\citep [option:text,text]
%TC:macro~\citet [option:text,text]
%TC:envir table 0 1
%TC:envir table* 0 1
%TC:envir tabular [ignore] word
%TC:envir displaymath 0 word
%TC:envir math 0 word
%TC:envir comment 0 0
%%
%% The first command in your LaTeX source must be the \documentclass
%% command.
%%
%% For submission and review of your manuscript please change the
%% command to \documentclass[manuscript, screen, review]{acmart}.
%%
%% When submitting camera ready or to TAPS, please change the command
%% to \documentclass[sigconf]{acmart} or whichever template is required
%% for your publication.
%%
%%
\documentclass[sigplan,10pt]{acmart}
\renewcommand\footnotetextcopyrightpermission[1]{}
\pagestyle{plain}
%%
%% \BibTeX command to typeset BibTeX logo in the docs
\AtBeginDocument{%
  \providecommand\BibTeX{{%
    Bib\TeX}}}

%% Rights management information.  This information is sent to you
%% when you complete the rights form.  These commands have SAMPLE
%% values in them; it is your responsibility as an author to replace
%% the commands and values with those provided to you when you
%% complete the rights form.
\setcopyright{acmlicensed}
\copyrightyear{2027}
\acmYear{2027}
\acmDOI{XXXXXXX.XXXXXXX}
\acmConference[EuroSys'27]{European Conference on Computer Systems}{April 19-23, 2027}{Rabat, Morocco}

%%
%% These commands are for a JOURNAL article.
%% \acmJournal{JDS}
%% \acmVolume{37}
%% \acmNumber{4}
%% \acmArticle{111}
%% \acmMonth{8}

%%
%% Submission ID.
%% Use this when submitting an article to a sponsored event. You'll
%% receive a unique submission ID from the organizers
%% of the event, and this ID should be used as the parameter to this command.
%%\acmSubmissionID{123-A56-BU3}
% \usepackage{showframe}
\usepackage{tabularx}
\usepackage{amsmath}
\usepackage{algorithm}
\usepackage{algorithmic}
\usepackage{subcaption}
\usepackage{stfloats}
\usepackage[normalem]{ulem}
\newtheorem{theorem}{Theorem}
\newtheorem{definition}{Definition}
\renewcommand{\qedsymbol}{$\blacksquare$}
\renewcommand{\algorithmicrequire}{ \textbf{Input:}}
\renewcommand{\algorithmicensure}{ \textbf{Output:}}
\renewcommand\thesubfigure{(\alph{subfigure})}
\captionsetup[sub]{labelformat=simple}

\bibliographystyle{plain}
%%
%% end of the preamble, start of the body of the document source.
\begin{document}

%%
%% The "title" command has an optional parameter,
%% allowing the author to define a "short title" to be used in page headers.
\title{Mosaic: Transparent Stitching of Stranded Resources across Multiple Hosts in Virtualization}

\author{Yanyan Li$^\dagger$, Zhiqian Xu$^\dagger$, Yawen Wang$^\dagger$, Hai Jiang$^{\dagger*}$, Keqin Li$^\diamondsuit$}
\affiliation{
    \institution{$^\dagger$Beijing University of Posts and Telecommunications \\
                 $^\diamondsuit$State University of New York}
}

%%
%% By default, the full list of authors will be used in the page
%% headers. Often, this list is too long, and will overlap
%% other information printed in the page headers. This command allows
%% the author to define a more concise list
%% of authors' names for this purpose.
%% \renewcommand{\shortauthors}{Trovato et al.}
%%
%% Article type: Research, Review, Discussion, Invited or position
\acmArticleType{Review}
%%
%% Links to code and data
%% \acmCodeLink{https://github.com/borisveytsman/acmart}
%% \acmDataLink{htps://zenodo.org/link}
%%
%% Authors' contribution
%% \acmContributions{BT and GKMT designed the study; LT, VB, and AP
%%  conducted the experiments, BR, HC, CP and JS analyzed the results,
%%  JPK developed analytical predictions, all authors participated in
%%  writing the manuscript.}
%%
%% Sometimes the addresses are too long to fit on the page.  In this
%% case uncomment the lines below and fill them accodingly.
%%
%% \authorsaddresses{Corresponding author: Ben Trovato,
%% \href{mailto:trovato@corporation.com}{trovato@corporation.com};
%% Institute for Clarity in Documentation, P.O. Box 1212, Dublin,
%% Ohio, USA, 43017-6221}
%%
%%
%% Keywords. The author(s) should pick words that accurately describe
%% the work being presented. Separate the keywords with commas.
\keywords{Cross-host, Resource Stitching, Virtualization}

\settopmatter{printfolios=true}

\begin{abstract}
In datacenters, the growing diversity of workloads causes some resources to turn into unallocatable stranded ones. Proactive schedulers and resource disaggregation have been designed to prevent or eliminate such stranding, but unpredictable agentic workloads and the high cost of disaggregation impede their deployment in production. This paper proposes Mosaic, which transparently stitches stranded resources across multiple hosts into an allocatable resource set. Mosaic is designed with a scheduling wrapper and a runtime system. The scheduling wrapper introduces a decomposition module to relax the indivisibility constraint of existing cluster schedulers, and the runtime system intercepts and forwards I/O requests at the para-virtualization emulator, requiring no modifications to either host or guest operating systems. We prototype Mosaic and evaluate it against a local KVM baseline and two kernel-based remote storage protocols (NVMe-oF and iSCSI). Mosaic retains 95\% of the local random-read throughput, converges with the kernel protocols under network RTT, and keeps the memory overhead below 10\%. On a TPCC macro-benchmark, Mosaic slows down by 16--44\% depending on concurrency and remains the fastest remote protocol at low concurrency.
\end{abstract}
\maketitle

\section{Introduction}\label{intro}
Datacenters have become the foundational infrastructure of digital society, supporting plenty of business~\cite{business1, business2} and computing demands~\cite{ai, cloud, hpc}. Boosting resource utilization of datacenters is a longstanding challenge across research and industry~\cite{mlaas, coach, legoos, feditor, protean}. Unused resources can be classified as unallocated, underutilized, and stranded~\cite{coach} ones. Unallocated resources are intentionally reserved to ensure service-level agreements~\cite{lazy, derm} and for cluster management~\cite{dynamo, mgr}. Underutilized resources are allocated to Virtual Machines (VMs) but not used~\cite{harvestvm, memory}. Stranded resources cannot be allocated since other resources in the same Physical Machine (PM) are fully allocated~\cite{coach, legoos, faasmem}. Solutions for unallocated and underutilized resources are deeply researched and stably deployed~\cite{lazy, derm, harvestvm}. However, resource stranding is still a common issue in datacenters, which has become more severe due to the growing diversity of workloads~\cite{agentcgroup, pond}.

Proactive schedulers~\cite{coach, resourcecentral, protean} and resource dis\-aggrega\-tion~\cite{legoos, pond, fractos} are two effective methods to alleviate resource stranding at the level of cluster scheduling and host operating system, respectively. However, proactive schedulers require resource usage patterns of workloads, while resource disaggregation demands massive hardware and software modifications. Both are difficult to be deployed. Mosaic, proposed in this paper, turns stranded resources into available ones by stitching them into an allocatable resource set in virtualization, as shown in Fig.~\ref{consolidation}. It is a reactive solution for stranded resources with transparency to both host and guest operating systems.

\begin{figure}[htbp]
    \centering
    \includegraphics[width=0.9\linewidth]{src/pic/consolidation.pdf}
    \caption{Resource stitching.}
    \label{consolidation}
\end{figure}

Mosaic exposes stitched resources through the para-virtualization framework, which natively separates the host and the guest and connects them via virtqueues. This separation lets Mosaic intercept the data path in the emulator and forward device accesses to another machine. We call the machine providing CPU cores the local side, and the machines providing peripherals the remote sides. As shown in Fig.~\ref{overall}, the expanded emulator consists of a local forwarding module and a remote forwarding module. The local module sends guest requests and handles remote responses; the remote module invokes devices and returns the results. To keep the crossing overhead low, Mosaic pipelines the forwarding path, batches small requests, applies zero-copy transmission above an adaptive size threshold, and pools workers so that multiple virtqueues proceed concurrently. An inflight window is further maintained so that outstanding requests hide the network latency.

\begin{figure}[htbp]
    \centering
    \includegraphics[width=0.9\linewidth]{src/pic/overall.pdf}
    \caption{The overall architecture of Mosaic.}
    \label{overall}
\end{figure}

We implement a prototype of Mosaic and evaluate it against a local KVM baseline and two kernel-based remote storage protocols, NVMe-oF and iSCSI. On microbenchmarks, Mosaic retains 95\% and 98\% of the local random-read and random-write throughput once the worker pool has at least two workers; under injected RTT up to 2\,ms, its throughput converges with NVMe-oF and iSCSI to the window/RTT bound, while the memory overhead stays below 10\%. On a TPCC macro-benchmark, Mosaic slows down by 16--44\% depending on concurrency yet remains the fastest remote protocol at low concurrency. On I/O-intensive Spark workloads, it stays within 8--16\% of KVM.

This paper makes the following contributions:
\begin{enumerate}
    \item Mosaic, a virtualization-layer resource stitching abstraction that allows cross-host resource usage, is proposed to turn stranded resources into allocatable ones.
    \item A cross-host forwarding architecture, based on para-virtualization framework, is designed, supporting pipe\-lined request processing.
    \item A low-overhead prototype of Mosaic is implemented and evaluated\footnote{anonymous code link}.
\end{enumerate}

This paper is structured as follows. Section~\ref{analysis} architecturally analyzes stranded resources. Section~\ref{runtimesystem} details the design of Mosaic. Section~\ref{eval} presents the evaluation results. Section~\ref{relatedwork} discusses related work, and Section~\ref{conclu} concludes.

\section{Analysis}\label{analysis}
\subsection{Stranded Resources in Different Managers}
Resource management in datacenters is completed via the coordination among the cluster scheduler, the host OS, and the emulator, as shown in Fig.~\ref{mapping}. The cluster scheduler assigns work nodes to VMs based on its scheduling objective, which is related to the states of Physical Machines (PMs) and VMs' resource demands. Each host OS manages fine-grained allocations on its PM. The emulator emulates virtual resources for the guest OS, with various dataplane acceleration mechanisms. Since emulators can only emulate allocated hardware, researchers have tried to solve the resource stranding issue in the cluster scheduler (i.e., proactive schedulers~\cite{autopilot, audible, proactive}) or the host OS (i.e., resource disaggregation~\cite{fractos, legoos, clio, pond}).

\begin{figure}[htbp]
    \centering
    \includegraphics[width=0.9\linewidth]{src/pic/mapping.pdf}
    \caption{Resource management in datacenters.}
    \label{mapping}
\end{figure}

The cluster scheduler traces the resource states of all PMs and schedules VM requests. The resource utilization of VMs and PMs can be modeled as time-series curves. A PM is excluded from the placement of a VM if the superposition of their utilization curves exceeds the allocation threshold. When a PM cannot satisfy any pending VM request, its remaining unallocated hardware goes unused and becomes stranded. To mitigate this, proactive schedulers group VMs with complementary resource demands into a single scheduling unit, smoothing out dynamic utilization fluctuations. By aligning the resource ratio of a unit with PM provisioning, resources are fully allocated simultaneously. However, as stated above, emerging unpredictable workloads lack stable statistical patterns, making curve abstraction impossible and rendering these proactive schedulers ineffective.

The host OS manages multiple types of hardware through different subsystems. Hardware availability can be viewed through the lens of these subsystems: when a resource is fully allocated, its corresponding subsystem rejects further requests, and a PM is excluded if a required subsystem rejects the VM's requests. When a PM cannot satisfy any pending VM request, resources whose subsystems still have capacity become stranded. Resource disaggregation disseminates resources and their OS subsystems to distributed PMs by decoupling hardware: although available resources are spread across multiple machines, their subsystems can be grouped to serve one VM, and since the usages of different resources are independent, stranding is essentially eliminated. On existing blade servers, disaggregation connects stranded resources by adding glue logic in the OS kernel; however, practical deployment is difficult because it requires intrusive kernel modifications and hardware-specific blade specifications.

\subsection{Design Principles}
Stitching a resource across hosts introduces two requirements beyond conventional VM I/O. First, the remote resource must remain transparent to the guest. Second, the additional cost of crossing the host boundary must be carefully controlled. We therefore follow three principles in designing Mosaic Runtime.

\textbf{P1: Transparency. }Mosaic should preserve the existing VM device abstraction. The guest continues to access a stitched resource through an unmodified interface, without knowing whether the underlying resource is local or remote. Also, the host is not aware of guest requests are forwarded to remote machines.

\textbf{P2: Low Crossing Overhead.} A remote I/O request incurs additional network, data-movement, synchronization, and buffer-management costs. Mosaic therefore pipelines request processing and amortizes per-request overhead through batching, while avoiding unnecessary data movement and buffer allocation.

\textbf{P3: Latency Tolerance.} Unlike local VirtIO, remote requests must tolerate network latency. Mosaic maintains sufficient outstanding requests and processes them concurrently so that network communication, remote device processing, and local request handling can overlap.

\section{Design and Implementation}\label{runtimesystem}
Based on the three principles, we choose para-virtualization as the base emulation model, since it is the only one that keeps I/O transparent to both sides. To reduce the overhead of crossing the host boundary, Mosaic applies zero-copy transmission, batching, and buffer pooling. For latency tolerance, Mosaic pipelines the forwarding path, pools workers to enable concurrency across virtqueues, and maintains an inflight window so that network latency overlaps with device processing.

\subsection{Why Para-Virtualization}
Para-virtualization (Virtio) is the only device-emulation model in the virtualization layer that keeps I/O transparent to both the host and the guest. Fig.~\ref{devicemodel} illustrates the three ways of emulating peripherals: full virtualization, pass-through, and para-virtualization. Full Virtualization (FV) emulates a complete device in software. Pass-Through (PT) exposes physical devices directly to virtual machines. Para-Virtualization (PV) introduces a cooperative interface between the guest and the emulator. Virtio offers a frontend-backend split abstraction: VM applications access devices through guest drivers, while backend drivers in the emulator invoke host drivers to handle I/O requests. This separation is exactly what Mosaic intercepts: by diverting requests at the emulator, the guest-visible device interface and the host-side drivers remain untouched, and the underlying cross-host forwarding is invisible to both operating systems.

\begin{figure}[htbp]
    \centering
    \includegraphics[width=0.9\linewidth]{src/pic/devicemodel.pdf}
    \caption{Device models in the QEMU emulator.}
    \label{devicemodel}
\end{figure}

\subsection{System Overview}
Mosaic exploits the separation between VirtIO front-end and back-end drivers to stitch resources across hosts at the virtualization layer. The front-end driver and the back-end driver of a stitched device execute on different machines and are connected through Mosaic's forwarding runtime. We refer to the machine hosting the guest CPU and QEMU as the local side, and the machine hosting the physical resource and its back-end driver as the remote side. From the guest's perspective, the remote resource remains an ordinary VirtIO device; Mosaic hides its physical location below the guest-visible device abstraction.

Figure~X shows the architecture of Mosaic. The expanded emulator consists of a local forwarding module and a remote forwarding module. The local module receives requests from VirtIO queues and forwards them to the remote side, while the remote module receives the requests and submits them to the physical device; device responses follow the reverse path. Since VirtIO devices commonly expose multiple virtqueues to support parallel I/O, Mosaic assigns an independent connection to each virtqueue, avoiding contention between otherwise independent request streams.

The design realizes the three principles of Section~\ref{analysis}. The existing VirtIO interface provides guest transparency. The forwarding pipeline, together with batching, zero-copy transmission, and buffer pooling, reduces the overhead of crossing the host boundary. Worker pooling and an inflight window maintain sufficient request concurrency to tolerate network latency. The mechanisms are described along the request path in the following sections.

% Using the separation of para-virtualization, Mosaic places front-end driver and back-end driver in different machines, and connects them with two forwarding modules. For clarity, the machine providing CPU cores to run guest is called as local side. All front-end drivers are executed on local side, while back-end drivers of stranded hardware runs on remote sides. Via this connection, stranded resources are stitched as an allocatable resource set. Additionally, Mosaic implements multiple mechanisms to satisfy design principle \textbf{2} and \textbf{3}.

% Many virtio devices uses multiple virtqueues to enable parallel device accessing. Therefore, Mosaic assigns each virtqueue with a socket to avoid contention among multiple virtqueues. Additionally, worker pooling mechanism is introduced to separate transactions of different virtqueues. Each arrived notifier is dispatched to a worker for sending requests or receiving responses.

% Request forwarding, including sending and receiving, is decoupled and asynchronized. Device requests flows between senders and receivers via lock-free buffers. The local forwarding module uses a Single Producer Single Consumer (SPSC) buffer to record requests that will generate responses. The remote forwarding module introduces a lock-free queue to record request that are arrived but not submit to devices.

% Requests could carry data with different sizes. To reduce redundant data copy when transmitting, Mosaic applies zero copy optimization for requests, whose data is larger than a threshold. This threshold is defined to avoid zero copy overhead outweighting the network performance. For requests, whose data is smaller than the threshold, Mosaic batches them to average the network latency. Via the two optimizations, all request transmission are accelerated.

% The remote forwarding module maintains a buffer pool to avoid overhead introduced by heap memory management. Since device submission is mutex, the receiver in remote forwarding batches requests, which arrives before finishing last batch, to continue network transmission, leaving backpressure away from local side.

\subsection{End-to-End Workflow}
The interaction workflow between the local QEMU and the remote stub is illustrated in Fig.~\ref{workflow}. This architecture expands the native Virtio ring buffer, which is confined to single-machine shared memory, into a cross-machine workflow.

\begin{figure}[htbp]
    \centering
    \includegraphics[width=\linewidth]{src/pic/workflow.pdf}
    \caption{The detailed workflow of Mosaic runtime.}
    \label{workflow}
\end{figure}

The virtio frontend driver in the guest sends an I/O request by writing the shared \texttt{vring} buffer of the corresponding \texttt{Virtqueue}, and then invokes a \texttt{kick} operation that issues an MMIO or PIO write, triggering a \texttt{VM\_exit}. The KVM subsystem checks the \texttt{exit\_reason} and notifies the emulator process \textcircled{1}. Since each \texttt{virtqueue} is shared between the guest and the emulator, backend drivers read I/O requests and write responses through the unified \texttt{VirtQueueElement} data structure. Upon sensing the notification, the overridden handler of Mosaic fetches the request as a \texttt{VirtQueueElement} through a framework interface \textcircled{2}. After analyzing the memory layout and other metadata, the local QEMU forwards the request by invoking a router module \textcircled{3}. The router is optimized with zero-copy and asynchronous pipelining to reduce communication overhead.

On the remote side, the stub parses the metadata in the payload and reconstructs the memory layout of the \texttt{VirtQueueElement}, avoiding a single large buffer. It then constructs a pseudo-local request and injects it into the backend driver \textcircled{4}. Unaware of the network translation, the backend driver processes the element and issues standard system calls to the remote host OS to operate the physical hardware \textcircled{5}. Once the physical I/O completes \textcircled{6}, the backend driver writes the resulting data back into the element and triggers a completion callback \textcircled{7}. Receiving this completion, the remote stub transmits the updated payload back to the local QEMU \textcircled{8}.

The router module in the local QEMU matches the \texttt{VirtQueueElement} against an ordered ring buffer, injects the response into the corresponding guest physical memory, and emits a completion signal into the local \texttt{virtqueue}. The guest driver is finally kicked through an \texttt{irq} \texttt{ioeventfd} \textcircled{9}.

\subsection{Local Forwarding}
The local forwarding module bridges the guest-facing VirtIO path and the cross-host transport. Its main challenge is to process requests concurrently without introducing additional synchronization and data-movement overhead into QEMU's I/O path.

\textbf{SPSC Pipelining and the Inflight Window.} Transactions for request forwarding and response handling are decoupled through a lock-free Single Producer Single Consumer (SPSC) buffer. A notification from the guest wakes a sending handler that retrieves available descriptors from the virtqueue and assembles them into network packages. If a device response is expected, the sender records a slot in the SPSC buffer; otherwise it kicks the guest directly. The response handler polls the socket and, because request flow is not ordered, traverses the sending window in the SPSC buffer to match an arrived response to its slot. Before kicking the guest, the receiver resorts the slots to remove holes. The SPSC buffer acts as an inflight window: Mosaic caps the number of outstanding requests at $I$, bounding memory usage and matching the concurrency that the device and the network can sustain.

\textbf{Worker Pool and Dispatching.} Stages of the forwarding pipeline are dispatched to threads taken from a worker pool. All remote devices of a Mosaic VM share the same pool of $W$ workers. A simple hash mapping assigns each virtqueue to a worker, avoiding idle polling and preventing hash-conflicting virtqueues from contending on the same handler.

\textbf{Batching and Adaptive Zero Copy.} Requests carry payloads of different sizes. For small requests, the fixed cost of transmitting each request dominates the actual data transfer, so Mosaic batches multiple small requests before transmission to amortize the per-request communication overhead. For requests with large payloads, Mosaic uses a zero-copy path to avoid unnecessary data movement, and retains the copy path for small payloads because the setup cost of zero-copy can outweigh its benefit; the threshold (1\,MB in our prototype) is chosen according to the relative cost of copying versus zero-copy setup. Additionally, releasing an SPSC slot requires a double check, since the kernel holds guest pages until the zero-copy completion signal arrives.

Together, SPSC pipelining, worker pooling, batching, and adaptive zero copy allow the local forwarding path to maintain high request throughput while limiting the additional cost introduced by remote forwarding.

\subsection{Remote Forwarding}
The remote forwarding module bridges the network-facing request stream and the physical device back-end. Unlike the local side, the remote module must also handle the synchronization and memory-management costs of submitting requests to the physical device.

\textbf{Pipelining, Worker Pool, and Zero Copy.} Similar to the local side, the receiving and sending handlers on the remote side are decoupled and pipelined. The receiver parses packages and reconstructs virtio-compliant device requests, while the sender forwards device responses. Both handlers are dispatched to workers from a shared worker pool with a simple hash mapping. The remote side does not batch response sending, delivering each response as soon as the device finishes processing it. Zero-copy transmission is used on the response path as well.

\textbf{Queueing and Batched Device Submission.} Although asynchronous backends allow multiple requests to be processed concurrently, submitting a request to the device takes a submission lock, so submission itself is serialized. Directly submitting each arriving request would expose this serialization to the network-facing path and bound the device concurrency by the arrival rate. Mosaic therefore maintains a submission queue with a limit $Q$ in the remote stub: requests that arrive while a previous submission is in progress are accumulated and submitted together when the current submission completes. $Q$ bounds the number of requests that can be concurrently in flight at the device, decoupling network reception from device submission, while batching the accumulated requests amortizes the submission-lock cost.

\textbf{Buffer Pool and Coalescing.} Mosaic maintains a buffer pool for request and response buffers. Without pooling, every forwarded request would trigger heap allocation and reclamation on the remote fast path; reusing preallocated buffers avoids this repeated memory-management overhead and provides stable buffer availability under high request rates. For requests whose payload is scattered across multiple guest pages, Mosaic optionally coalesces the segments into a single contiguous buffer, enabling one large device I/O and a single response iovec; this removes the heap churn of large-block forwarding and improves latency for blocks of 256\,KB and above.

The response path follows the reverse direction. After the physical device completes a request, Mosaic places the completion into the corresponding response path and forwards it through the per-virtqueue connection to the local side. The local forwarding module then completes the corresponding VirtIO descriptor and returns control to QEMU.

\subsection{Implementation Summary}
We implement the Mosaic runtime system by extending QEMU's virtio with around 4K lines of code. It requires zero modifications to the guest OS drivers, the native device backend logic, and the KVM hardware acceleration path. We first hook the control plane by registering QOM properties on the \texttt{VirtioDeviceClass}; the setters for these properties dynamically override the native call-back functions of the \texttt{ioeventfd} handlers to process Mosaic transactions. For data-plane interception, a \texttt{remote\_ctx} pointer is injected into the core structure \texttt{VirtQueue}. Eight interfaces of the virtio framework are modified, seamlessly diverting local I/O processing to remote handlers. To minimize Mosaic's interference with other local devices, all related threads are registered as QEMU aio iothreads, so Mosaic's execution never blocks other transactions in the main loop. On the remote side, we leverage an extremely trimmed machine type \texttt{x-remote-machine} that contains only a PCI bus to plug stranded hardware pieces.


\section{Evaluation}\label{eval}
This section presents the performance evaluation of Mosaic scheduling wrapper and runtime system, using micro- and macro-benchmarks. We ran all experiments on a machine with Intel Core i5-13600 K CPU (5.10 GHz) with 32 GB memory and 500 GB NVMe SSD. Network latencies are deterministically injected by Linux traffic control.

\subsection{Experimental Setup}
The Mosaic runtime is evaluated with a 4-vCPU guest and a remote NVMe device exposed through four virtqueues. All I/O accesses the raw physical device with O\_DIRECT, and host page caches are dropped before each run to exclude caching effects. Because the underlying NVMe exhibits run-to-run bandwidth variance of up to $\pm$55\%, we always run comparisons back-to-back, anchor each configuration with a host-level direct-I/O measurement, and report medians of three runs. Mosaic is compared against three baselines: native KVM with a local virtio-blk device, NVMe-oF over TCP, and iSCSI, the latter two serving the same physical device through kernel initiator stacks.

Microbenchmarks use fio 3.41 with libaio, O\_DIRECT, and 60-second steady-state windows (\texttt{time\_based}, single job unless stated). The Mosaic configuration space covers the worker-pool size $W \in \{1,2,4\}$, the inflight window $I$, the remote queue limit $Q$, a zero-copy threshold of 1\,MB, and a batch size of 16 requests.

\subsection{Microbenchmarks}

\textbf{Throughput versus the local baseline.} We first run four 4\,KB workloads (random read, random write, sequential read, sequential write) at depth 16 with four jobs, varying only the worker-pool size $W$; the KVM baseline serves the same device locally. With $W{=}1$, all four workloads degrade substantially: a single worker serializes the four virtqueues, and random read drops to 153.7K IOPS (77\% of the 199.0K baseline), random write to 40.3K (42\% of 96.0K), sequential read to 89.0K (40\% of 220.0K), and sequential write to 67.8K (56\% of 122.0K). Increasing $W$ to 2 already restores most performance: random read and random write reach 175.9K and 73.5K IOPS, sequential read 135.9K, and sequential write 87.0K. At $W{=}4$, random read and random write reach 189.5K and 94.2K, i.e., 95\% and 98\% of the local baseline, and sequential write 96.6K (79\%); the three random-friendly workloads are statistically within the run-to-run variance of the baseline. Sequential read is the only workload that retains a real loss of 148.3K versus 220K IOPS ($-$33\%) even at $W{=}4$: its latency is dominated by the per-byte transfer cost of the remote path, which does not amortize with concurrency. Random write is the only workload that improves monotonically with $W$ (40.3$\rightarrow$73.5$\rightarrow$94.2K), reflecting the increasing parallelism of device submission.

% ===== FIGURE PLACEHOLDER: m2g_grouped_iops.svg (IOPS by workload, W=1/2/4 vs KVM) =====

\textbf{Device concurrency and the inflight window.} Mosaic bounds the number of requests submitted to the physical device with a queue limit $Q$ in the remote stub. We sweep $Q \in \{1,16,32,64\}$ for guest depths of 1, 16, 32, and 64 under 4\,KB random read. At depth 1, all settings deliver the identical 5.5K IOPS (the 178--180\,$\mu$s serial baseline). With $Q{=}1$, the stub issues strictly one asynchronous I/O at a time and throughput is capped at 7.0--7.2K regardless of depth, since all requests serialize on the 140\,$\mu$s device latency. Raising $Q$ to 16 unlocks device parallelism and yields a 10$\times$ jump to 55.6K (depth 16), rising to 66.7K and 71.2K at depths 32 and 64. $Q{=}32$ behaves like $Q{=}16$ (55.4/77.0/82.8K). The inflight window $I$ provides additional queueing slack that hides RTT: keeping $Q{=}64$ and guest depth 64, enlarging the window from 64 to 128 raises throughput from 65.9K to 92.2K IOPS ($+$40\%), reaching 91\% of the KVM baseline (101.4K); without this slack the window is exactly full and every network round trip is exposed. The effective throughput is thus bounded by the weakest of device concurrency $\min(\mathrm{inflight},Q)$ and protocol concurrency $\min(\mathrm{inflight},I)$: $Q$ controls how much parallelism reaches the device, while $I$ controls how much queueing slack hides network latency.

% ===== FIGURE PLACEHOLDER: m6_iodepth_curves.svg (IOPS vs depth, per Q) =====

\textbf{Latency tolerance under network RTT.} We inject one-way delays between 0 and 950\,$\mu$s (RTT 0.1--2\,ms) on the loopback interface and run a 70/30 read/write mix at depth 32 for window sizes 16 and 32, and for the two kernel protocols. At RTT 0.1\,ms, Mosaic with window 32 sustains 53.5K IOPS, on par with NVMe-oF (48.2K) and 82\% of KVM (65.1K); iSCSI reaches 61.4K. As RTT grows, all remote protocols converge to the same throughput curve governed by window/RTT: at 2\,ms, Mosaic, NVMe-oF, and iSCSI deliver 14.1K, 15.1K, and 13.8K IOPS respectively, close to the theoretical bound of 32 requests / 2\,ms $=16$K (efficiency 88--94\%). A 20$\times$ RTT increase therefore costs only 3.2--4.4$\times$ throughput, entirely explained by the Little's-law bound; the local KVM baseline stays flat at 65.1K since it never crosses the network. Halving the window halves throughput across the entire RTT range (ratio 0.40--0.54, e.g., 29.2K vs.\ 53.5K at 0.1\,ms and 5.8K vs.\ 14.1K at 2\,ms), confirming that the inflight window is the first-order bottleneck of the forwarding path: it limits concurrency at low RTT and outstanding requests at high RTT.

% ===== FIGURE PLACEHOLDER: m3_rtt_iops.svg (IOPS vs RTT, 4 curves) =====

\textbf{Memory overhead.} We measure the memory footprint with a single in-flight request to isolate steady state, for block sizes of 64K, 512K, 1M, and 2M. Mosaic adds only 1--6\% RSS to the local emulator over the corresponding KVM guest (e.g., 5--9\% total at the peak), and the remote stub maintains a constant 34--48\,MB across all sizes; zero-copy kernel buffers (skmem) remain at the kilobyte level. The overhead is independent of request size and does not grow with concurrency. Interestingly, the zero-copy path is not more memory-intensive than the copy path at low concurrency, since both must hold the same number of in-flight buffers; the benefit of zero copy appears only by eliminating per-request data movement, not by reducing buffering.

% ===== FIGURE PLACEHOLDER: m1mem_bar_all.svg (memory by bs, local/stub/KVM) =====

\textbf{Large requests and buffer pooling.} For large blocks, where per-byte cost dominates, we run a synchronous workload at block sizes from 16K to 2M. KVM's per-request latency grows 4.5$\times$ (19.8 to 88.5\,$\mu$s), while Mosaic without buffer pooling grows 16.8$\times$ (48.7 to 818\,$\mu$s) and with buffer pooling 13.5$\times$ (49.8 to 671\,$\mu$s): the remote path multiplies the disk-intrinsic cost by data movement, transfer, and submission overhead. The bandwidth ratio of Mosaic to KVM therefore shrinks from 40\% at 16K to 13\% at 2M. Buffer pooling recovers most of this loss: at 256K, 1M, and 2M it reduces latency by 10\%, 26\%, and 18\% respectively (172.3$\rightarrow$154.5, 506.1$\rightarrow$375.6, 818.0$\rightarrow$671.3\,$\mu$s), while at 64K and below the difference is within noise. Reusing preallocated buffers avoids the heap churn (per-request allocation and reclamation) that dominates large-block forwarding; requests below 256K are too short for this overhead to matter.

% ===== FIGURE PLACEHOLDER: m5_pdiff.svg (P=on benefit vs bs) =====

\subsection{Macrobenchmarks}
We evaluate a sysbench TPCC workload (20-second warm-up followed by 60-second measurement, page caches dropped per run) with 1, 2, and 4 client threads; the guest data disk is served through each protocol in turn, back-to-back on the same day.

% ===== FIGURE PLACEHOLDER: macro_txn_abs.svg (txns/s by threads, 4 protocols) =====

At one thread, Mosaic achieves 185.6 txns/s, 23\% below the KVM baseline (240.0) but higher than NVMe-oF (158.3) and iSCSI (109.3): with little concurrency the fixed per-transaction overhead dominates, and Mosaic's lightweight virtio forwarding is cheaper than the kernel protocol stacks. At two threads all protocols scale (Mosaic 252.2, NVMe-oF 221.5, iSCSI 171.6 vs.\ KVM 298.7), and Mosaic remains the fastest remote path. As concurrency grows, the inflight window becomes the bottleneck: at four threads Mosaic delivers 248.5 txns/s ($-$44\% vs.\ KVM 446.3), comparable to iSCSI (249.5) and slightly below NVMe-oF (273.8), and all three remote paths plateau because the guest-side window caps the outstanding transactions.

% ===== FIGURE PLACEHOLDER: macro_p95_abs.svg (p95 by threads, 4 protocols) =====

Tail latency degrades more than throughput. Mosaic's p95 transaction latency rises from 10.65\,ms at one thread (+31\% vs.\ KVM's 8.13\,ms) to 44.17\,ms at four threads (+113\% vs.\ 20.74\,ms), while its average latency grows from 5.39\,ms (+29\%) to 16.09\,ms (+80\%). NVMe-oF shows p95 of 13.22/20.74/37.56\,ms and iSCSI 30.26/38.94/45.79\,ms across the three thread counts; iSCSI's p95 reaches 3.7$\times$ KVM at one thread. The degradation ranking of tail latency (Mosaic $<$ NVMe-oF $<$ iSCSI at low concurrency) mirrors the ranking of fixed per-transaction cost, confirming that the p95 inflation comes from the same forwarding path.

\subsection{End-to-End Applications}
We run three data-intensive applications on the data disk: PageRank on a 1M-page Wikipedia graph (10 iterations), ALS matrix factorization over a movie-rating dataset (20 training iterations), and XGBoost over a 6M-row synthetic dataset (16 rounds). For ALS and XGBoost we force each iteration to re-read the dataset from the data disk and drop the page cache, keeping I/O on the critical path; the workloads run with 1, 2, and 4 Spark cores.

% ===== FIGURE PLACEHOLDER: app_pagerank_cores.svg (PageRank total time by cores) =====

With a 150\,$\mu$s RTT, PageRank shows no protocol difference: total times at 1/2/4 cores are 39.2/24.9/21.7\,s (Mosaic), 39.6/27.3/23.0\,s (KVM), 38.2/24.3/20.9\,s (NVMe-oF), and 38.3/24.6/21.1\,s (iSCSI). Computation dominates and I/O never becomes the bottleneck, so the 150\,$\mu$s RTT is fully hidden.

% ===== FIGURE PLACEHOLDER: app_als_cores.svg (ALS total, stacked build+train) =====

When I/O is forced back into the critical path, the protocols separate. For ALS, Mosaic stays within 8.0--13.8\% of KVM across core counts (79.7/69.7/65.8\,s vs.\ 73.8/62.1/57.8\,s), and its four-core build phase is actually the fastest among all protocols (9.4\,s), because four parallel readers saturate the remote disk; the loss concentrates in the serial training phase, which re-reads the dataset from disk each iteration. For XGBoost, Mosaic loses 8.4--15.9\% (60.0/57.7/57.2\,s vs.\ 55.4/50.6/49.4\,s), with the gap again concentrated in the rescan-driven training.

% ===== FIGURE PLACEHOLDER: app_xgb_cores.svg (xgb total, stacked build+train) =====

iSCSI is the worst in both: 75--150\% slower, because its single-queue backend saturates under parallel re-reads. These results confirm that, once device parallelism is unlocked, Mosaic's forwarding overhead remains small relative to kernel remote-storage protocols even for I/O-intensive applications, and that its remaining loss scales with the volume of data crossing the remote path rather than with core count.

\section{Related Work}\label{relatedwork}
\subsection{Virtualization-based Forwarding}
Virtualization systems have explored separating the virtual device interface from the process that implements the device. QEMU Multi-process moves device emulation out of the main QEMU process and communicates with the device process through a Unix domain socket, improving isolation while retaining the conventional virtual device interface \cite{qemu-multiprocess}. VFIO-user similarly defines an interface between a VMM and a userspace device implementation, allowing a device to be implemented outside the VMM process \cite{vfio-user}. These systems demonstrate that the virtualization layer can serve as a boundary between a VM and the implementation of its devices.

Mosaic builds on this separation principle but extends it across physical hosts. Instead of merely moving device execution to another process on the same machine, Mosaic places the VirtIO front-end and back-end on different hosts and forwards VirtIO transactions over the network. This enables a physical device that is stranded on one host to serve a VM running on another host without modifying the guest-visible device interface. The resulting cross-host forwarding path, however, introduces network latency and additional software overhead that do not exist in local process-level separation. Mosaic therefore focuses on designing an efficient forwarding runtime using asynchronous pipelining, per-virtqueue isolation, batching, and adaptive zero-copy.

\subsection{Hardware-assisted Disaggregation}
Hardware-assisted disaggregation decouples resources from individual hosts using specialized interconnects and resource fabrics. Recent systems extend this approach beyond memory to physical devices. Oasis pools PCIe devices, including NICs and SSDs, across hosts using CXL infrastructure \cite{oasis}, while DxPU disaggregates GPUs from their physical hosts and exposes them through a shared GPU pool \cite{dxpu}. These systems provide efficient remote access by moving resource sharing into the hardware or fabric layer.

Mosaic takes a complementary software-based approach. Remote devices remain attached to their original hosts, while the virtualization layer exposes them to VMs through the existing VirtIO abstraction. Mosaic therefore requires neither a dedicated device fabric nor specialized hardware, making it applicable to conventional virtualized clusters. The tradeoff is the additional network and software overhead of forwarding device requests, which Mosaic addresses through its optimized runtime.

\subsection{Software-based Resource Disaggregation}
Resource disaggregation also exists purely in software, especially for storage and memory. Kernel protocols such as NVMe-oF~\cite{nvmeof} and iSCSI~\cite{iscsi} assign stranded storage by encapsulating raw I/O commands into network protocols, and are widely deployed. For memory, RDMA-based swapping, e.g., InfiniSwap~\cite{infiniswap} and FastSwap~\cite{fastswap}, uses remote memory as a swap region, while AIFM~\cite{aifm} adds aggressive prefetching and program analysis to reduce swap traffic. LegoOS~\cite{legoos} goes further with a split-kernel design that disseminates computing, memory, and storage across distributed machines and allocates each resource on demand. These systems share a common cost: they either change the guest-visible interface or require modifications to the OS kernel. Mosaic instead reuses the unmodified VirtIO stack and the host-side device drivers: it forwards I/O at the emulator, leaving both the guest and the host untouched, at the price of the forwarding overhead quantified in Section~\ref{eval}.

\section{Discussion and Conclusion}\label{conclu}
\noindent\textbf{Discussion.} Mosaic relies on the standardized frontend-backend separation of the virtio ecosystem. While virtio has become the de facto standard for mainstream datacenter I/O~\cite{virtio}, its penetration into specialized hardware remains uneven. Specifically, for emerging AI accelerators~\cite{nvidia, google, ascend}, hardware vendors are highly incentivized by commercial interests to maintain vertically integrated software stacks rather than developing open virtio drivers. Limited vendor support for standardized virtio interfaces restricts Mosaic's deployment on some closed-source AI workloads.

\noindent\textbf{Conclusion.} In this paper, we present Mosaic, a reactive solution to stranded resources. It transparently stitches stranded resources across multiple hosts into an allocatable resource set through the para-virtualization layer. A scheduling wrapper relaxes the indivisibility constraint of cluster schedulers, and a runtime system forwards device I/O across hosts without modifying guest or host operating systems. Evaluations show that Mosaic retains 95\% of the local random-read throughput, converges with kernel remote-storage protocols under network RTT, and slows a TPCC macro-benchmark by 16--44\% depending on concurrency, while remaining the fastest remote protocol at low concurrency. The remaining loss concentrates at high concurrency, where the inflight window caps outstanding requests; improving this behavior, e.g., by integrating kernel-level data-plane acceleration, is left as future work.

\newpage

\bibliography{src/ref}

\end{document}
\endinput
%%
%% End of file sample-acmcp.tex'.
